\section{Introduction}

Fake news can impose costs well beyond an isolated inaccurate belief. In democratic settings, disinformation can distort public debate, intensify polarisation, and erode trust in elections, institutions, and journalism; the OECD therefore treats information integrity as an enabling condition for democratic participation and accountable government \citep{oecd2024}. Public-health decisions provide a second concrete example. In a randomised experiment in the United Kingdom and the United States, exposure to COVID-19 vaccine misinformation reduced the share of respondents who said they would definitely accept vaccination by 6.2 and 6.4 percentage points, respectively \citep{loomba2021}. These examples do not imply that every false publication changes behaviour or determines an election. They show why the reach and repetition of deceptive content can create societal consequences that justify careful, rights-aware intervention.

Online social networks allow news sources to reach large audiences with low publication and redistribution costs. The same infrastructure supports legitimate journalism, uncertain information, deliberately deceptive content, and mixtures of these forms. False news has been shown to diffuse farther and faster than true news in large Twitter cascades \citep{vosoughi2018}, while exposure and sharing are concentrated among subsets of users \citep{grinberg2019,guess2019}. These findings motivate interventions, but they do not by themselves reveal how a malicious source may adapt its observable presentation when platforms or users become more selective.

This paper focuses on strategies controlled by a news source that repeatedly disseminates false publications. Such a source may imitate attributes associated with established media: emotional engagement, political or social proximity, apparent information quality, or source credibility. We distinguish these perceived attributes from the objective probability that a publication is false. This separation is essential. A camouflage strategy can make a source more attractive without making its output more accurate. We use \textit{false news} for factually false news-like publications, \textit{misinformation} without assuming intent, and \textit{disinformation} only when deliberate deception is established. Because the model specifies behaviour rather than an internal intention state, ``malicious source'' is shorthand for the experimental source and not an empirical diagnosis of motive.

Observed social-media data are indispensable but make controlled counterfactuals difficult. Agent-based modelling (ABM) complements observation by representing users with heterogeneous states and network positions, repeated decisions and feedback, and by changing one policy or strategy while holding the remaining design constant \citep{macal2010,grimm2020}. Our model therefore serves as a computational laboratory for source strategies and mitigation policies; it is not intended to forecast X at the individual-user level.

We address four research questions:

\begin{description}
\item[RQ1.] How do source-imitation strategies affect malicious-source reach and fake-news dissemination? Platforms and users observe many source characteristics, but it is unclear which groups of characteristics are most exploitable or whether a source gaining audience necessarily increases network-wide false-news exposure. RQ1 contributes a controlled ranking of individual and combined imitation strategies while holding objective fake-news probability fixed.
\item[RQ2.] How do campaign timing and duration affect efficacy and persistence? Deceptive presentation may be adopted early, late, or temporarily, while detection and intervention occur after different delays. A full-run average confounds start time with exposure duration, so RQ2 contributes matched treatment and post-removal windows that distinguish opportunity during a campaign from persistence after it ends.
\item[RQ3.] How effectively does limiting experimental-source reach mitigate dissemination, and how do recommendation policies modify this relationship? Reach and recommendation are platform-controllable mechanisms, but interpersonal diffusion could complement or bypass direct visibility restrictions. An immediate and perfectly truth-aware mechanism is an informative boundary, not a realistic platform description. RQ3 therefore contributes a joint assessment of graded reach limits, discovery-only recommendation, asymmetric flagging, delayed and imperfect labels, and recommendation magnitude.
\item[RQ4.] Are strategy conclusions robust to parameter and structural uncertainty? Memory persistence and contact counts are uncertain behavioural assumptions rather than direct policy levers, and conclusions may also depend on source behaviour, choice scaling, user activity and network construction. RQ4 contributes local, global and structural sensitivity analyses that distinguish stable conclusions from effects requiring stronger behavioural qualifications.
\end{description}

The contribution is threefold. First, we provide a modular ABM in which named source strategies selectively copy observable attributes while publication veracity remains unchanged. Second, we use common random-number seeds in a staged computational experiment to compare strategy, timing, reach and recommendation policies and to screen parameter and structural uncertainty. Third, we distinguish four outcomes that policy discussions can otherwise conflate: false reposts per agent-period, false reposts among actual decisions, experimental-source share among decisions, and participation.

The remainder of the paper is organised as follows. Section~2 reviews fake-news dissemination, ABM, and related simulation applications. Section~3 describes the model and experimental design using the ODD protocol. Section~4 reports the four research-question analyses. Section~5 documents software verification and conceptual and behavioural validation. Section~6 develops societal and policy implications while delimiting the claims supported by the experiments. Section~7 concludes and identifies the next validation steps.
